\bmtchapitre{Glossaire}{Relier un terme à son sens précis.}{Repères de vocabulaire}
\begin{description}[style=nextline,leftmargin=1em]
\item[Antécédent]Élément $x$ du domaine tel que $f(x)=y$; un même $y$ peut en avoir plusieurs.
\item[Asymptote]Droite dont l’écart à la courbe vérifie la limite donnée au chapitre~\ref{t11s-c03}.
\item[Barycentre]Point unique annulant une somme de vecteurs pondérés, lorsque la somme des poids est non nulle.
\item[Bijection]Application injective et surjective.
\item[Contre-exemple]Cas appartenant au domaine et réfutant une proposition universelle.
\item[Contraposée]Pour $P\Rightarrow Q$, l’implication $\neg Q\Rightarrow\neg P$, logiquement équivalente.
\item[Dérivée]Fonction des nombres dérivés, sur les points où ils existent.
\item[Équiprobabilité]Égalité des probabilités de toutes les issues du modèle fini considéré.
\item[Invariant]Propriété conservée à chaque tour d’un algorithme.
\item[Médiane]Valeur centrale de la série ordonnée selon la convention précisée.
\item[Monotone]Croissante ou décroissante sur l’intervalle, ou sur les rangs, précisés.
\item[Paramètre]Valeur fixée pendant la résolution, dont on étudie ensuite les différents cas.
\item[Point fixe]Point dont l’image par une transformation est lui-même.
\item[Produit vectoriel]Vecteur orthogonal aux deux facteurs; son sens dépend de l’ordre et de l’orientation du repère.
\item[Réciproque]Pour $P\Rightarrow Q$, l’implication $Q\Rightarrow P$, dont la vérité est à établir séparément.
\item[Variance]Moyenne des carrés des écarts à la moyenne, exprimée en unités au carré.
\end{description}
